% !TEX root = CoeneDylanBP.tex
%%=============================================================================
%% Objectdetectie en landingsframe detectie
%%=============================================================================

\section{Object- en landingsframe detectie}
\label{sec:yolo}

Met een gekalibreerde camera en berekende homografiematrix is het doel van deze fase om een \textit{first proof-of-concept} te realiseren voor de meting van horizontale verplaatsing (HD). Hiervoor wordt een objectdetectiemodel ingezet om de springer per frame te lokaliseren en zijn relatieve verplaatsing doorheen de sprong te bepalen. Daarnaast wordt geprobeerd om het landingsframe te detecteren aan de hand van de verticale snelheid van het zwaartepunt van de bounding box.

Als detectiemodel wordt YOLOv11 gebruikt~\autocite{Jocher_Ultralytics_YOLO_2023}. Dit model is geschikt voor dit project vanwege de hoge verwerkingssnelheid en de bewezen nauwkeurigheid bij het herkennen van personen in beweging. De volledige pipeline verloopt als volgt:

\begin{enumerate}
    \item YOLOv11 detecteert de springer in elk frame en geeft een bounding box terug.
    \item Het zwaartepunt (middelpunt) van de bounding box wordt berekend in pixelcoördinaten $(u, v)$.
    \item Doorheen de frames wordt de verticale snelheid van het zwaartepunt berekend.
    \item Het moment van landing wordt geschat door te kijken naar het moment dat deze snelheid plotseling afneemt.
\end{enumerate}

Het gebruik van een bounding box zwaartepunt als referentiepunt om de horizontale verplaatsing te bepalen kent een inherente beperking: de grootte en positie van de bounding box vari{\"e}ren afhankelijk van de houding van de springer (bv.\ gestrekte armen vs.\ een gekromde sprong). Dit kan leiden tot een systematische fout in de geprojecteerde positie. Deze beperking geeft een logische aanleiding om in Fase~\ref{sec:pose} over te stappen naar Pose Estimation.

De gemeten landingszones worden vergeleken met manueel bepaalde referentieposities. De nauwkeurigheid wordt uitgedrukt als een percentage correct geclassificeerde landingszones (accuracy) en als een gemiddelde afwijking in meters. % TODO: Dit is nog niet zeker
